%%%%%%%%%%%%%%%%%%%%%%%%%%%%%%%%%%%%%%%%%%%%%%%%%%%%%%%%%%%%%%%%%%%%%%%%%%%%%%%
% Derivation Attempt v1: 5D→4D Newton Constant
% Actual calculation, not roadmap
%
% Status: WORKING / OPEN
%%%%%%%%%%%%%%%%%%%%%%%%%%%%%%%%%%%%%%%%%%%%%%%%%%%%%%%%%%%%%%%%%%%%%%%%%%%%%%%
\documentclass[11pt,a4paper]{article}

\usepackage[utf8]{inputenc}
\usepackage[T1]{fontenc}
\usepackage{amsmath,amssymb,amsthm}
\usepackage{geometry}
\geometry{margin=2.5cm}
\usepackage{hyperref}
\hypersetup{colorlinks=true, linkcolor=blue!60!black, citecolor=blue!60!black}
\usepackage{booktabs}
\usepackage{enumitem}
\usepackage{tcolorbox}
\tcbuselibrary{skins,breakable}

% Epistemic tags
\newcommand{\tagBL}{\textcolor{blue!70!black}{[BL]}}
\newcommand{\tagI}{\textcolor{green!50!black}{[I]}}
\newcommand{\tagP}{\textcolor{orange!80!black}{[P]}}
\newcommand{\tagOPEN}{\textcolor{red!70!black}{[OPEN]}}
\newcommand{\tagDc}{\textcolor{purple!70!black}{[Dc]}}

% Box environments
\newtcolorbox{statusbox}{
  colback=yellow!5!white,
  colframe=yellow!50!black,
  fonttitle=\bfseries,
  title=Epistemic Status,
  breakable
}

\newtcolorbox{resultbox}{
  colback=green!5!white,
  colframe=green!50!black,
  fonttitle=\bfseries,
  breakable
}

\newtcolorbox{failbox}{
  colback=red!5!white,
  colframe=red!50!black,
  fonttitle=\bfseries,
  title=Derivation Incomplete,
  breakable
}

\title{\textbf{Derivation Attempt v1:}\\[0.3em]
5D $\to$ 4D Newton Constant\\[0.5em]
\large Actual Calculation (BLOCK-003)}
\author{EDC Research Program}
\date{2026-02-02\\[0.3em]\small\textit{Status: WORKING / OPEN}}

\begin{document}

\maketitle

%===============================================================================
\begin{abstract}
%===============================================================================
This document attempts an explicit derivation of the effective 4D Newton constant $G_N$ from the 5D bulk action used in EDC.
We linearize around a flat background, apply Israel junction conditions at the brane, solve for the static weak-field potential, and extract $G_N$.
The derivation succeeds in obtaining a standard brane-world result: $G_N = \kappa_5^2 / (8\pi L)$ for a compact extra dimension of size $L$, or $G_N \sim \kappa_5^4 \sigma$ for RS-type warped geometries.
However, the EDC-specific scaling involving powers 12/13 does \emph{not} emerge from this standard calculation.
We identify exactly where additional EDC-specific input would be required, marking BLOCK-003 as still \tagOPEN{}.
\end{abstract}

%===============================================================================
\section{Setup and Notation}
\label{sec:setup}
%===============================================================================

\begin{statusbox}
\textbf{Status: \tagBL{}/\tagI{}.}
Standard brane-world setup with EDC coordinate conventions.
\end{statusbox}

\paragraph{Coordinates and signature \tagBL{}.}
We use a 5D bulk manifold $\mathcal{M}^5$ with coordinates
\begin{equation}
    X^A = (x^\mu, \xi), \qquad A = 0,1,2,3,5, \qquad \mu = 0,1,2,3,
\end{equation}
where $x^\mu = (t, x^i)$ are 4D coordinates and $\xi$ is the extra-dimensional coordinate.
Signature convention: $(-,+,+,+,+)$.

\paragraph{Brane location \tagI{}.}
A 3-brane $\Sigma$ is located at $\xi = 0$.
Brane observers experience the induced metric $\gamma_{\mu\nu} = g_{\mu\nu}|_{\xi=0}$.

\paragraph{5D gravitational coupling \tagBL{}.}
The fundamental coupling is $\kappa_5^2 = 8\pi G_5$, where $G_5$ has dimensions $[\text{length}]^3$.

\paragraph{Action \tagBL{}.}
The total action is:
\begin{equation}
    S = S_{\text{bulk}} + S_{\text{GHY}} + S_{\text{brane}},
    \label{eq:total_action}
\end{equation}
with
\begin{align}
    S_{\text{bulk}} &= \frac{1}{2\kappa_5^2} \int_{\mathcal{M}^5} d^5X \sqrt{-g} \left( R^{(5)} - 2\Lambda_5 \right), \label{eq:S_bulk} \\
    S_{\text{GHY}} &= \frac{1}{\kappa_5^2} \oint_{\partial\mathcal{M}} d^4x \sqrt{|h|} \, K, \label{eq:S_GHY} \\
    S_{\text{brane}} &= -\sigma \int_\Sigma d^4x \sqrt{-\gamma} + S_{\text{matter}}[\gamma, \psi]. \label{eq:S_brane}
\end{align}
Here $\sigma$ is the brane tension and $K = h^{ab} K_{ab}$ is the trace of the extrinsic curvature \cite{York1972,GH1977}.

\paragraph{Newtonian limit definition \tagBL{}.}
The ``Newtonian limit'' means: static, weak-field, non-relativistic source.
On the brane, we seek a potential $\Phi(r)$ such that the 4D metric is
\begin{equation}
    ds^2_{\text{brane}} \approx -(1 + 2\Phi)\,dt^2 + (1 - 2\Phi)\,d\vec{x}^2,
\end{equation}
with $|\Phi| \ll 1$ and $\Phi(r) \to -G_N M / r$ at large $r$ for a point mass $M$.

%===============================================================================
\section{Background and Linearization}
\label{sec:linearize}
%===============================================================================

\begin{statusbox}
\textbf{Status: \tagP{}/\tagBL{}.}
We choose a flat background (no warp) with $\Lambda_5 = 0$ for simplicity.
This is a \emph{postulate} that limits generality.
\end{statusbox}

\paragraph{Background choice \tagP{}.}
We take the simplest case: flat 5D Minkowski bulk with $\Lambda_5 = 0$:
\begin{equation}
    \bar{g}_{AB} = \eta_{AB} = \text{diag}(-1, +1, +1, +1, +1).
    \label{eq:background}
\end{equation}
The brane sits at $\xi = 0$ and we assume $\mathbb{Z}_2$ orbifold symmetry $\xi \to -\xi$.
The brane tension $\sigma$ is set to zero for the background (or treated as a first-order perturbation).

\paragraph{Perturbation \tagBL{}.}
Write
\begin{equation}
    g_{AB} = \eta_{AB} + h_{AB}, \qquad |h_{AB}| \ll 1.
    \label{eq:perturbation}
\end{equation}

\paragraph{Gauge choice \tagBL{}.}
We use the transverse-traceless (TT) gauge adapted to the brane:
\begin{equation}
    \partial^A h_{AB} = \frac{1}{2} \partial_B h, \qquad h \equiv \eta^{AB} h_{AB}.
    \label{eq:gauge}
\end{equation}
This is the 5D analog of de Donder (harmonic) gauge.

\paragraph{Linearized Einstein equations \tagDc{}.}
The 5D Einstein tensor linearized around flat space is \cite{MTW,Wald}:
\begin{equation}
    G_{AB}^{(1)} = -\frac{1}{2} \Box_5 \bar{h}_{AB},
    \label{eq:linearized_G}
\end{equation}
where $\bar{h}_{AB} = h_{AB} - \frac{1}{2}\eta_{AB} h$ is the trace-reversed perturbation and
\begin{equation}
    \Box_5 = \eta^{CD}\partial_C \partial_D = -\partial_t^2 + \nabla^2_3 + \partial_\xi^2.
\end{equation}
In our gauge, the linearized vacuum field equations become:
\begin{equation}
    \boxed{\Box_5 \bar{h}_{AB} = 0 \quad (\text{bulk, away from brane}).}
    \label{eq:bulk_eom}
\end{equation}

%===============================================================================
\section{Israel Junction Conditions}
\label{sec:junction}
%===============================================================================

\begin{statusbox}
\textbf{Status: \tagBL{}.}
Standard Israel junction conditions \cite{Israel1966}.
\end{statusbox}

\paragraph{Extrinsic curvature \tagBL{}.}
For a hypersurface at $\xi = 0$ with unit normal $n^A = (0,0,0,0,1)$, the extrinsic curvature is:
\begin{equation}
    K_{\mu\nu} = -\frac{1}{2} \partial_\xi g_{\mu\nu} \big|_{\xi=0}.
    \label{eq:extrinsic}
\end{equation}
With $\mathbb{Z}_2$ symmetry, the jump across the brane is $[K_{\mu\nu}] = 2 K_{\mu\nu}^+$ where $K^+$ is evaluated from $\xi > 0$.

\paragraph{Israel junction conditions \tagBL{}.}
The discontinuity in extrinsic curvature is related to the brane stress-energy \cite{Israel1966}:
\begin{equation}
    [K_{\mu\nu}] - [K] \gamma_{\mu\nu} = -\kappa_5^2 S_{\mu\nu},
    \label{eq:israel}
\end{equation}
where $S_{\mu\nu}$ is the brane stress-energy tensor.
For a brane with tension $\sigma$ and localized matter $\tau_{\mu\nu}$:
\begin{equation}
    S_{\mu\nu} = -\sigma \gamma_{\mu\nu} + \tau_{\mu\nu}.
\end{equation}
Substituting:
\begin{equation}
    [K_{\mu\nu}] - [K] \gamma_{\mu\nu} = \kappa_5^2 \sigma \gamma_{\mu\nu} - \kappa_5^2 \tau_{\mu\nu}.
    \label{eq:israel_expanded}
\end{equation}

\paragraph{Linearized junction conditions \tagDc{}.}
At linear order with $\mathbb{Z}_2$ symmetry, using \eqref{eq:extrinsic}:
\begin{equation}
    \partial_\xi h_{\mu\nu}\big|_{\xi=0^+} = -\kappa_5^2 \left( \tau_{\mu\nu} - \frac{1}{3} \tau \eta_{\mu\nu} \right),
    \label{eq:junction_linear}
\end{equation}
where $\tau = \eta^{\mu\nu} \tau_{\mu\nu}$ and we have used that for a $\mathbb{Z}_2$ brane, $[K_{\mu\nu}] = 2 K_{\mu\nu}^+ = -\partial_\xi h_{\mu\nu}|_{0^+}$.

%===============================================================================
\section{Static Potential on the Brane}
\label{sec:potential}
%===============================================================================

\begin{statusbox}
\textbf{Status: \tagDc{}/\tagOPEN{}.}
We derive the Green's function structure but encounter a divergence requiring regularization.
\end{statusbox}

\paragraph{Point mass source \tagBL{}.}
Place a static point mass $M$ at the origin on the brane:
\begin{equation}
    \tau_{00} = M \delta^{(3)}(\vec{x}), \qquad \tau_{ij} = 0, \qquad \tau = -M \delta^{(3)}(\vec{x}).
    \label{eq:source}
\end{equation}

\paragraph{Mode expansion \tagDc{}.}
For the static case ($\partial_t = 0$), the bulk equation \eqref{eq:bulk_eom} becomes:
\begin{equation}
    \left( \nabla^2_3 + \partial_\xi^2 \right) \bar{h}_{\mu\nu} = 0.
    \label{eq:static_bulk}
\end{equation}
Fourier transform in the 3D spatial coordinates:
\begin{equation}
    \bar{h}_{\mu\nu}(\vec{x}, \xi) = \int \frac{d^3k}{(2\pi)^3} \, \tilde{h}_{\mu\nu}(\vec{k}, \xi) \, e^{i\vec{k}\cdot\vec{x}}.
\end{equation}
The bulk equation becomes:
\begin{equation}
    \left( \partial_\xi^2 - k^2 \right) \tilde{h}_{\mu\nu} = 0, \qquad k = |\vec{k}|.
    \label{eq:bulk_fourier}
\end{equation}
Solution with $\mathbb{Z}_2$ symmetry and decay as $|\xi| \to \infty$:
\begin{equation}
    \tilde{h}_{\mu\nu}(\vec{k}, \xi) = \tilde{h}_{\mu\nu}^{(0)}(\vec{k}) \, e^{-k|\xi|}.
    \label{eq:bulk_solution}
\end{equation}

\paragraph{Apply junction condition \tagDc{}.}
From \eqref{eq:junction_linear}, the derivative jump at $\xi = 0$ gives:
\begin{equation}
    \partial_\xi \tilde{h}_{\mu\nu}\big|_{0^+} = -k \, \tilde{h}_{\mu\nu}^{(0)} = -\kappa_5^2 \left( \tilde{\tau}_{\mu\nu} - \frac{1}{3} \tilde{\tau} \eta_{\mu\nu} \right).
\end{equation}
Thus:
\begin{equation}
    \tilde{h}_{\mu\nu}^{(0)}(\vec{k}) = \frac{\kappa_5^2}{k} \left( \tilde{\tau}_{\mu\nu} - \frac{1}{3} \tilde{\tau} \eta_{\mu\nu} \right).
    \label{eq:h_brane}
\end{equation}

\paragraph{Brane potential \tagDc{}.}
For the $00$-component with the point source \eqref{eq:source}:
\begin{equation}
    \tilde{\tau}_{00} = M, \qquad \tilde{\tau} = -M.
\end{equation}
Hence:
\begin{equation}
    \tilde{h}_{00}^{(0)} = \frac{\kappa_5^2}{k} \left( M + \frac{M}{3} \right) = \frac{4\kappa_5^2 M}{3k}.
\end{equation}
Inverting the Fourier transform:
\begin{equation}
    h_{00}(\vec{x}, 0) = \frac{4\kappa_5^2 M}{3} \int \frac{d^3k}{(2\pi)^3} \frac{e^{i\vec{k}\cdot\vec{x}}}{k}.
    \label{eq:h00_integral}
\end{equation}

\begin{failbox}
The integral \eqref{eq:h00_integral} is:
\begin{equation}
    \int \frac{d^3k}{(2\pi)^3} \frac{e^{i\vec{k}\cdot\vec{x}}}{k} = \frac{1}{2\pi^2 r^2}.
    \label{eq:integral_result}
\end{equation}
This gives $h_{00} \sim 1/r^2$, \textbf{not} the Newtonian $1/r$ behavior!

\textbf{Missing element:} In a non-compact extra dimension, gravity ``leaks'' into the bulk, producing $1/r^2$ at all scales. To recover 4D $1/r$ gravity, we need either:
\begin{enumerate}[nosep]
    \item A \textbf{compact} extra dimension of size $L$ (Kaluza--Klein), or
    \item A \textbf{warped} geometry that localizes the graviton zero mode (Randall--Sundrum).
\end{enumerate}
\end{failbox}

%===============================================================================
\section{Compact Extra Dimension: Recovery of 4D Gravity}
\label{sec:compact}
%===============================================================================

\begin{statusbox}
\textbf{Status: \tagP{}/\tagDc{}.}
We \emph{assume} a compact $\xi$-direction of size $L$. This is a postulate.
\end{statusbox}

\paragraph{Compactification \tagP{}.}
Let $\xi \in [0, L]$ with $\mathbb{Z}_2$ orbifold at both ends (or periodic with period $2L$).
The bulk solution now involves a discrete mode sum:
\begin{equation}
    \tilde{h}_{\mu\nu}(\vec{k}, \xi) = \sum_{n=0}^\infty \tilde{h}_{\mu\nu}^{(n)}(\vec{k}) \cos\left(\frac{n\pi\xi}{L}\right).
\end{equation}

\paragraph{Zero mode dominance \tagDc{}.}
At distances $r \gg L$, only the $n=0$ (zero) mode contributes significantly.
The zero mode is constant in $\xi$ and satisfies:
\begin{equation}
    \nabla^2_3 \bar{h}_{\mu\nu}^{(0)} = -\frac{\kappa_5^2}{L} \left( \tau_{\mu\nu} - \frac{1}{3}\tau\eta_{\mu\nu} \right) \delta^{(3)}(\vec{x}),
\end{equation}
where the factor $1/L$ comes from integrating over the extra dimension.

\paragraph{4D Poisson equation \tagDc{}.}
This is exactly the 4D linearized Einstein equation:
\begin{equation}
    \nabla^2 h_{00}^{(0)} = -\frac{4\kappa_5^2 M}{3L} \delta^{(3)}(\vec{x}).
\end{equation}
Solution:
\begin{equation}
    h_{00}^{(0)}(r) = \frac{\kappa_5^2 M}{3\pi L} \cdot \frac{1}{r}.
    \label{eq:h00_4D}
\end{equation}

\paragraph{Identification of $G_N$ \tagDc{}.}
Comparing with the Newtonian potential $\Phi = -G_N M / r$ and using $h_{00} = -2\Phi$:
\begin{equation}
    -2\Phi = \frac{\kappa_5^2 M}{3\pi L r} \quad \Rightarrow \quad \Phi = -\frac{\kappa_5^2 M}{6\pi L r}.
\end{equation}
Thus:
\begin{equation}
    \boxed{G_N = \frac{\kappa_5^2}{6\pi L} = \frac{G_5}{L}.}
    \label{eq:GN_result}
\end{equation}

\begin{resultbox}
\textbf{Result (Kaluza--Klein):}
\begin{equation}
    G_N = \frac{\kappa_5^2}{6\pi L} = \frac{4G_5}{3L},
\end{equation}
where $L$ is the compactification length.

\textbf{Epistemic status:} \tagDc{} — derived from linearized 5D GR with compact extra dimension.

\textbf{Required input:} The length scale $L$ must be specified independently.
\end{resultbox}

%===============================================================================
\section{Warped Geometry Alternative}
\label{sec:warped}
%===============================================================================

\begin{statusbox}
\textbf{Status: \tagP{}/\tagBL{}.}
Standard RS result quoted; full derivation requires solving warped bulk equations.
\end{statusbox}

\paragraph{Randall--Sundrum setup \tagBL{}.}
For a warped metric $ds^2 = e^{-2A(\xi)} \eta_{\mu\nu} dx^\mu dx^\nu + d\xi^2$ with AdS$_5$ bulk \cite{RS1999}:
\begin{equation}
    G_N = \frac{\kappa_5^4 \sigma}{48\pi} \quad \text{(RS2 limit)},
    \label{eq:GN_RS}
\end{equation}
where the effective 4D Planck mass is set by the warp factor localization.

\paragraph{Key difference from EDC \tagI{}.}
The RS result \eqref{eq:GN_RS} gives $G_N \sim \kappa_5^4 \sigma$.
This is \emph{not} the structure appearing in EDC gravitational claims, which involve specific power combinations (12, 13) of other parameters.

%===============================================================================
\section{Confrontation with EDC Scaling}
\label{sec:edc_scaling}
%===============================================================================

\begin{failbox}
\textbf{The derivation does NOT produce EDC-specific powers (12/13).}

From standard brane-world physics, we obtain:
\begin{itemize}[nosep]
    \item \textbf{Kaluza--Klein:} $G_N = \kappa_5^2 / (6\pi L)$ — depends on compactification scale $L$.
    \item \textbf{Randall--Sundrum:} $G_N \sim \kappa_5^4 \sigma$ — depends on tension and 5D coupling.
\end{itemize}

Neither of these has the form claimed in EDC gravitational sector discussions.
To obtain powers 12/13 would require:
\begin{enumerate}
    \item A specific EDC-derived relationship between $L$ (or warp scale) and other EDC parameters ($\sigma$, $R_\xi$, membrane geometry), \textbf{not derived here} \tagOPEN{}.
    \item Or a fundamentally different mechanism beyond standard linearized brane-world gravity.
\end{enumerate}

\textbf{Conclusion:} BLOCK-003 remains \tagOPEN{}. The power-law structure is \textbf{not derived} by this calculation.
\end{failbox}

%===============================================================================
\section{Epistemic Ledger}
\label{sec:ledger}
%===============================================================================

\begin{center}
\renewcommand{\arraystretch}{1.3}
\begin{tabular}{p{4cm}p{1.5cm}p{7cm}}
\toprule
\textbf{Step} & \textbf{Status} & \textbf{Notes} \\
\midrule
5D action & \tagBL{} & Standard Einstein--Hilbert + GHY + brane \\
Flat background choice & \tagP{} & Simplification; limits generality \\
Linearized equations & \tagDc{} & Derived from action \\
Israel junction conditions & \tagBL{} & Standard result \cite{Israel1966} \\
Mode expansion & \tagDc{} & Standard Fourier methods \\
Non-compact: $1/r^2$ & \tagDc{} & Derived; incorrect for 4D gravity \\
Compactification $L$ & \tagP{} & Postulated to recover $1/r$ \\
$G_N = \kappa_5^2/(6\pi L)$ & \tagDc{} & Derived given compactification \\
RS warped result & \tagBL{} & Quoted from literature \\
EDC powers 12/13 & \tagOPEN{} & Not derived; no mechanism identified \\
\bottomrule
\end{tabular}
\end{center}

%===============================================================================
\section{What Is Required to Close BLOCK-003}
\label{sec:required}
%===============================================================================

\paragraph{Missing lemmas \tagOPEN{}.}
\begin{enumerate}
    \item \textbf{Length scale origin:} The compactification length $L$ (or warp scale in RS) must be derived from EDC-internal parameters, not imposed externally.

    \item \textbf{Power structure:} If EDC claims $G_N \sim \sigma^a R_\xi^b$ with specific powers $a, b$, then either:
    \begin{itemize}[nosep]
        \item A derivation of $L = f(\sigma, R_\xi, \ldots)$ from EDC geometry, or
        \item A non-standard gravitational mechanism specific to EDC.
    \end{itemize}

    \item \textbf{Anti-circularity:} The derivation must not use $G_N$ as input.
    Currently, \eqref{eq:GN_result} satisfies this (uses only $\kappa_5$, $L$), but $L$ itself may hide circularity if defined via 4D observations.
\end{enumerate}

%===============================================================================
\section{Separation of Concerns}
\label{sec:separation}
%===============================================================================

This derivation attempt (v1) is a \emph{working document} linked to the frozen BLOCK-003 program note in \texttt{paper\_gravity\_block003/}.

\begin{itemize}[nosep]
    \item The program note defines acceptance criteria AC-G1 through AC-G10.
    \item This document addresses AC-G1 (action), AC-G2 (junction), AC-G3 (background), AC-G4 (linearization), AC-G6 (propagator structure), and partially AC-G7 ($G_N$ extraction).
    \item AC-G5 (KK modes), AC-G8 (numerical consistency), AC-G9 (parameter counting), AC-G10 (reproducibility) remain open or partially addressed.
\end{itemize}

The nuclear pinning monograph (\texttt{edc\_book\_2/}) is \textbf{not modified} by this work.

%===============================================================================
\section{Conclusion}
\label{sec:conclusion}
%===============================================================================

\begin{enumerate}
    \item We derived the standard brane-world result: $G_N = \kappa_5^2 / (6\pi L)$ for a compact extra dimension.
    \item The derivation shows explicitly why a non-compact extra dimension fails (gives $1/r^2$).
    \item The EDC-specific power structure (12/13) does \textbf{not emerge} from this standard calculation.
    \item BLOCK-003 remains \tagOPEN{}: closing it requires deriving the compactification/warp scale from EDC-internal parameters.
\end{enumerate}

%===============================================================================
% References
%===============================================================================

\begin{thebibliography}{9}

\bibitem{MTW}
C.~W. Misner, K.~S. Thorne, and J.~A. Wheeler,
\textit{Gravitation} (W.~H. Freeman, 1973).

\bibitem{Wald}
R.~M. Wald,
\textit{General Relativity} (University of Chicago Press, 1984).

\bibitem{York1972}
J.~W. York,
\textit{Phys.\ Rev.\ Lett.}\ \textbf{28}, 1082 (1972).

\bibitem{GH1977}
G.~W. Gibbons and S.~W. Hawking,
\textit{Phys.\ Rev.\ D}\ \textbf{15}, 2752 (1977).

\bibitem{Israel1966}
W.~Israel,
\textit{Nuovo Cimento B}\ \textbf{44}, 1 (1966).

\bibitem{RS1999}
L.~Randall and R.~Sundrum,
\textit{Phys.\ Rev.\ Lett.}\ \textbf{83}, 3370 (1999).

\end{thebibliography}

\end{document}
